\documentclass[11pt]{article}
\usepackage[a4paper,margin=25mm]{geometry}
\usepackage{amsmath,amssymb,amsthm,booktabs}
\newtheorem{theorem}{Theorem}\newtheorem{lemma}[theorem]{Lemma}
\newcommand{\Aut}{\operatorname{Aut}}\newcommand{\Cor}{\operatorname{Cor}}
\title{Common incidences determine eleven full geometry groups}
\author{}\date{}
\begin{document}\maketitle
Let $G=\mathrm{HS}$, and let $A=\Aut(G)$, in the acquired faithful
models of orders $44\,352\,000$ and $88\,704\,000$.
We use the reviewed classification and source coordinates for the
coset geometries in D.~Leemans, \emph{The residually weakly primitive
geometries of HS}, Australasian Journal of Combinatorics 33 (2005),
231--236, with the retained supplementary diagrams.
The original results remain attributed to that source.

\begin{theorem}
The following are the full type-preserving automorphism groups and
full correlation groups, respectively. Correlations may permute types.
\begin{center}
\begin{tabular}{ccll}\toprule
Representative&Source label&$\Aut(\Gamma)$&$\Cor(\Gamma)$\\\midrule
11&3.25&$A$&$A$\\
12&3.26&$A$&$A\times C_2$\\
23&3.14&$G$&$G$\\
26&3.11&$A$&$A\times C_2$\\
27&3.10&$A$&$A\times C_2$\\
29&3.7&$A$&$A$\\
49&4.32&$A$&$A$\\
50&4.33&$A$&$A$\\
55&4.15&$G$&$G\times C_2$\\
71&4.7&$A$&$A\times C_2$\\
74&4.3&$A$&$A$\\\bottomrule
\end{tabular}
\end{center}
The finite proof uses exact permutation arithmetic and the complete
graph automorphism algorithm specified below.
\end{theorem}

\paragraph{Common-incidence relations.}
Fix any type $a$ of size $n=1100$, with base object $b$ and stabilizer
$P_a$. The acquired group is transitive on that type. Let $O_k$ be
the orbits of $P_a$ on its objects, of sizes $d_k$; omit the singleton
$\{b\}$. They parametrize the orbits on ordered distinct point pairs,
of respective sizes $nd_k$. For another type $j$, let
$S_j=b^{P_j}$ be the point shadow of its base object. Define $e_{jk}$
to be the number of ordered distinct pairs in $S_j$ belonging to pair
orbit $k$. Then the number of type-$j$ objects incident with both
members of a pair in orbit $k$ is
\begin{equation}\label{counts}
                  m_{jk}=\frac{[G:P_j]e_{jk}}{nd_k}.
\end{equation}
Indeed, count triples consisting of a type-$j$ object and an ordered
pair in its shadow. Counting first by objects gives the numerator;
transitivity on the pair orbit makes the count per pair constant.
This argument counts objects even when two objects have identical shadows.
For $j=a$, the shadow convention is equality within that type,
so all off-diagonal counts are zero.

For a distinct pair, form the sorted multiset of pairs
$([G:P_j],m_{jk})$ over all types $j$. This signature is intrinsic:
a correlation merely permutes its entries and sends $a$ to a type of
the same size. Equal signatures partition the distinct point pairs
into intrinsic relations. In representatives 11 and 23 these relations
have valencies 280 and 819. In each remaining representative, at
every 1100-object anchor, there is exactly one relation of valency $105$.
Use the valency $280$ relation in the first two cases and the valency $105$
relation in the others. The literal checks verify symmetry, so each
chosen relation is an undirected graph.

\paragraph{Computing the relations in actual coordinates.}
The saved permutations on 1100 points define a bijective homomorphism
from the actual $G$ onto its image $E$. The specified base point has
stabilizer exactly the image of $P_a$. Recorded elements $t_x\in G$
send that base point to every point $x$, in order. These facts prove
the coset identification $P_ag\mapsto b^g$; the certificate need not
assume that a producer's enumeration was a correct transversal.
For each ordered pair $(x,y)$ in $S_j$, its orbit index is that of
$y^{t_x^{-1}}$ under $E_b$. Thus finite counts of these indices compute
$e_{jk}$ and hence~\eqref{counts}. Transporting the selected base
neighbour set by the $t_x$ gives every adjacency row. All specified
generators are checked on the entire graph.

\paragraph{Faithfulness of point restriction.}
We use the elementary shadow lemma from the reviewed common theory:
an equivariant encoding of $P_j\backslash G$ is injective precisely
when its base value has stabilizer $P_j$. Begin by encoding an object
by its point shadow. If an already identified type $k$ is needed,
adjoin the set of encodings of its type-$k$ neighbours. At the base
object of type $j$ this is exactly the $P_j$-orbit of the base
type-$k$ encoding. Its cardinality is checked to be
$[P_j:P_j\cap P_k]$.

The positive verifier reads the saved encodings in dependency order.
It admits an added neighbour set only after type $k$ has been proved
injectively encoded, verifies that entire orbit, and intersects the
corresponding set stabilizers. The final stabilizer must equal the
actual image of $P_j$. There are sixteen possible 1100-object anchors
in these eleven families; every type is identified at every anchor.
In particular, representative $49$ at its third type needs successive
neighbour sets to identify types $1$ and $2$. No direct-shadow assertion is
substituted for these iterations.
An automorphism fixing the point type pointwise therefore fixes every
object, by induction through the same dependencies. Restriction of
the full typed automorphism group to any one anchor is faithful.

\paragraph{The complete upper bound.}
For each of the sixteen graphs, a complete nauty computation returns
order $88\,704\,000$. Every returned generator is checked on every
adjacency row. In the literal permutation group $D$ they generate,
the verifier checks $|D|=88\,704\,000$, $E\leq D$ and $D'=E$.
Thus $E$ is characteristic in the full graph group. This last equality,
rather than agreement of orders alone, supplies the needed invariant
copy of the specified group.

Any correlation sends an anchor graph to the graph at its target
anchor, by the intrinsic common-incidence definition. If $c$ is such
a correlation and $g\in G$, characteristicity implies that the action
of $cgc^{-1}$ on the target anchor equals that of some $h\in G$.
The typed automorphism $h^{-1}cgc^{-1}$ fixes that anchor pointwise.
Faithfulness proves it is the identity. Consequently every full
correlation normalizes the specified $G$.

\paragraph{Exact normalizers and structures.}
Put $B=\bigcap_iP_i$, and let $S,T$ be the ordered and setwise
stabilizers of the parabolic family in the acquired $\Aut(G)$.
The reviewed coset normalizer theorem now applies to the full groups:
\[
 |\Aut(\Gamma)|=|G:B||S|,\qquad
 |\Cor(\Gamma)|=|G:B||T|.
\]
The same exhaustive transporter calculation tests every type permutation.
It verifies actual family images, not just subgroup orders.
For the representatives in table order the triples $(|B|,|S|,|T|)$ are
\[
\begin{gathered}
(16,32,32),\ (16,32,64),\ (72,72,72),\ (192,384,768),\\
(192,384,768),\ (192,384,384),\ (4,8,8),\ (4,8,8),\\
(8,8,16),\ (32,64,128),\ (64,128,128).
\end{gathered}
\]
In every row $S\cap G=B$. The centralizer of $G$ in the full
correlation group is the actual quotient $(T\cap G)/B$. Its order
is two exactly for representatives $12,26,27,55,71$, and is otherwise
one. In the former cases it is central and meets the typed group
trivially; the typed group and this centralizer have full product
order. Hence the displayed factors $C_2$ are direct factors.
The faithful conjugation action identifies the remaining factor with
$G$ or $\Aut(G)$ as indicated. This proves all group structures.

\paragraph{Attribution and verification boundary.}
The coset, shadow and normalizer arguments are reused from the common
theory. Formula~\eqref{counts} generalizes the previously reviewed
100-point count; sorted signatures allow arbitrary type interchanges.
The present contribution is this application and its explicit finite
identifications. The original classification, exact source-coordinate
joins, GAP's finite group algorithms and nauty's completeness are
acquired premises. Literal small-group controls compare the formula
with direct common-neighbour counts, including repeated shadows.
No new classification census, historical novelty, conjecture resolution,
Lean verification or whole-manuscript completion is asserted here.
\end{document}
